% BHU PYQ -- Manually transcribed & error-corrected
% Paper: MTB-501 : Mathematical Analysis
% Exam: B.A./B.Sc. (Hons) Semester V Examination, 2022-23

\documentclass[11pt,a4paper]{article}
\usepackage[a4paper,top=2.5cm,bottom=2.5cm,left=2.8cm,right=2.8cm,headheight=14pt]{geometry}
\usepackage{amsmath,amssymb}
\usepackage[T1]{fontenc}
\usepackage[utf8]{inputenc}
\usepackage{lmodern}
\usepackage{microtype}
\usepackage{fancyhdr}
\usepackage{enumitem}
\usepackage{hyperref}
\usepackage{lastpage}
\usepackage{parskip}

\hypersetup{colorlinks=true,urlcolor=black,linkcolor=black,
  pdftitle={MTB-501: Mathematical Analysis -- BHU B.A./B.Sc. Sem V 2022-23}}
\pagestyle{fancy}\fancyhf{}
\renewcommand{\headrulewidth}{0.4pt}\renewcommand{\footrulewidth}{0.4pt}
\fancyhead[L]{\small\textit{Banaras Hindu University}}
\fancyhead[R]{\small\textit{MTB-501: Mathematical Analysis}}
\fancyfoot[C]{\small Page \thepage\ of \pageref{LastPage}}

\newlist{parts}{enumerate}{2}
\setlist[parts,1]{label=(\alph*),leftmargin=*,itemsep=5pt,topsep=3pt}
\setlist[parts,2]{label=(\roman*),leftmargin=*,itemsep=3pt,topsep=2pt}
\newcommand{\pts}[1]{\hfill\textit{[#1]}}

\begin{document}
\begin{center}
  {\large\bfseries BANARAS HINDU UNIVERSITY}\\[2pt]
  {\normalsize B.A./B.Sc. (Hons) --- Semester V Examination, 2022-23}\\[6pt]
  \rule{0.8\linewidth}{0.5pt}\\[6pt]
  {\large\bfseries Paper No. MTB-501: Mathematical Analysis}\\[4pt]
  {\normalsize Time: 3 Hours \hfill Full Marks: 70}
\end{center}
\rule{\linewidth}{0.4pt}
\smallskip
\noindent\textit{Instructions: Answer \textbf{five} questions including Question~1 which is compulsory. Figures in right-hand margin indicate marks.}
\bigskip

\noindent\textbf{Question 1.} \textit{(Compulsory.)}
\begin{parts}
  \item Show that the limit of a convergent sequence is unique. \pts{2}
  \item Show that the sequence $(x_n)_{n=1}^\infty$, where $x_n = \left(1 + \dfrac{1}{n}\right)^n$, is monotonically increasing. \pts{2}
  \item Determine the values of $x$ for which the series $\sum_{n=0}^\infty \dfrac{x^n}{n!}$ is absolutely convergent. \pts{2}
  \item Test the convergence of the series $\sum_{n=1}^\infty \dfrac{\cos(1/n)}{n^2}$. \pts{2}
  \item Let $f(x)$ be defined by:
  \[
  f(x) = \begin{cases} \dfrac{\sin x}{x}, & x \neq 0 \\ 1, & x = 0 \end{cases}
  \]
  Prove that $f$ is Riemann integrable on $\mathbb{R}$ (on any closed bounded interval). \pts{2}
  \item Let $f(x) = x$ on $[0,1]$. If $P = \left\{0, \dfrac{1}{n}, \dfrac{2}{n}, \dots, 1\right\}$ is a partition of $[0,1]$, then find: (i)~$L(P, f)$, (ii)~$U(P, f)$, (iii)~$\underline{\int_0^1} f\,dx$, (iv)~$\overline{\int_0^1} f\,dx$. \pts{4}
  \item If $f(x)$ and $g(x)$ are Riemann integrable on $[a, b]$ such that $f(x) \le g(x)$ for all $x \in [a, b]$, show that $\int_a^b f\,dx \le \int_a^b g\,dx$. \pts{2}
  \item Test the convergence of $\int_0^1 \dfrac{x}{1-x}\,dx$. \pts{2}
  \item Test the convergence of $\int_0^\infty e^{-x^2}\,dx$. \pts{2}
  \item Show that $\int_0^\infty \dfrac{\sin^2 x}{x^2}\,dx = \dfrac{\pi}{2}$. \pts{2}
\end{parts}

\medskip\rule{\linewidth}{0.2pt}

\medskip\noindent\textbf{Question 2.}
\begin{parts}
  \item Prove that the necessary and sufficient condition for the convergence of a monotonic sequence of real numbers is that it is bounded. \pts{6}
  \item Prove that if a sequence $(x_n)_{n=1}^\infty$ of positive numbers converges to a limit $l$, then:
  \[
  \lim_{n\to\infty} (x_1 x_2 \cdots x_n)^{1/n} = l
  \]
  Deduce that $\lim_{n\to\infty} n^{1/n} = 1$. \pts{6}
\end{parts}

\medskip\rule{\linewidth}{0.2pt}

\medskip\noindent\textbf{Question 3.}
\begin{parts}
  \item State and prove Cauchy's general principle for the convergence of sequences of real numbers. \pts{6}
  \item Let $\sum_{n=1}^\infty a_n$ be a positive term series and let $\lim_{n\to\infty} (a_n)^{1/n} = l$. Then the series is:
  \begin{parts}
    \item convergent if $l < 1$;
    \item divergent if $l > 1$;
    \item the test fails if $l = 1$.
  \end{parts}
  Prove this result (Cauchy's root test). \pts{6}
\end{parts}

\medskip\rule{\linewidth}{0.2pt}

\medskip\noindent\textbf{Question 4.}
\begin{parts}
  \item Test the convergence of the series $\sum_{n=2}^\infty \dfrac{1}{n(\log n)^p}$. \pts{6}
  \item Define absolute and conditional convergence of a series. Show that the series $\sum_{n=1}^\infty \dfrac{(-1)^{n+1}}{n^p}$ is absolutely convergent for $p > 1$, but only conditionally convergent for $0 < p \le 1$. \pts{6}
\end{parts}

\medskip\rule{\linewidth}{0.2pt}

\medskip\noindent\textbf{Question 5.}
\begin{parts}
  \item Prove that a necessary and sufficient condition for the Riemann integrability of a bounded function $f$ is that for every $\varepsilon > 0$, there exists $\delta > 0$ such that for every partition $P$ of $[a, b]$ with norm $\|P\| < \delta$:
  \[
  U(P, f) - L(P, f) < \varepsilon
  \] \pts{6}
  \item Prove that every monotonic function on $[a, b]$ is Riemann integrable on $[a, b]$. Give an example of a Riemann integrable function on $[a, b]$ which is not monotonic on $[a, b]$. \pts{6}
\end{parts}

\medskip\rule{\linewidth}{0.2pt}

\medskip\noindent\textbf{Question 6.}
\begin{parts}
  \item Show that the function defined by:
  \[
  f(x) = \frac{1}{n+1} \quad \text{for } \frac{1}{n+1} \le x < \frac{1}{n},\ n = 1, 2, 3, \dots
  \]
  is Riemann integrable on $[0, 1]$. Also show that $\int_0^1 f(x)\,dx = \dfrac{\pi^2}{6} - 1$. \pts{6}
  \item If $\int_a^b f(x)\,dx$ and $\int_a^b g(x)\,dx$ both exist and $g(x)$ keeps the same sign throughout the interval of integration, then there exists a number $\mu$ lying between the bounds of $f(x)$ such that:
  \[
  \int_a^b f(x)g(x)\,dx = \mu \int_a^b g(x)\,dx
  \]
  (State and prove the Second Mean Value Theorem for integrals.) \pts{6}
\end{parts}

\medskip\rule{\linewidth}{0.2pt}

\medskip\noindent\textbf{Question 7.}
\begin{parts}
  \item State and prove Dirichlet's test for improper integrals. \pts{6}
  \item State and prove the Limit Comparison Test (LCT) for improper integrals. Examine the convergence of $\int_0^\infty \dfrac{\sin^p x}{x}\,dx$ for $p > 1$. \pts{6}
\end{parts}

\vfill
\rule{\linewidth}{0.4pt}
\begin{center}\small
  \href{https://bhu-exam.vercel.app/}{Click here to visit our website}\\[4pt]
  \href{https://me-aryan.vercel.app/}{Click here to visit our contributor}\\[4pt]
  \href{https://quashan-me.vercel.app/}{Click here to visit our contributor}
\end{center}

\end{document}
